% E14: 三準位(211・322・544)の chi_SR の f 依存(mu = 0.71e-12 eV)。データは calc/e14_three_level.py、e12_review.py
\begin{tikzpicture}
\begin{axis}[note, width=0.72\linewidth, height=0.42\linewidth, xmode=log, xmin=1e11, xmax=1e19, ymin=0.15, ymax=0.82,
  xlabel={$f\ [\mathrm{GeV}]$}, ylabel={年齢 $T$ のあとのスピン}, unbounded coords=discard,
  legend style={font=\scriptsize, at={(0.99,0.97)}, anchor=north east, fill=white, fill opacity=0.85, text opacity=1, draw=none}, legend cell align=left]
  \fill[cM!12] (axis cs:1e17,0.15) rectangle (axis cs:1e19,0.82);
  \node[font=\scriptsize, color=ink2] at (axis cs:1e18,0.33) {未確認};
  \addplot[ink2, line width=0.8pt, dashed] table[x=f, y=chi_GW241011_T9] {figs/data/e12_scan.dat}; \addlegendentry{E12(二準位)、GW241011 の主星}
  \addplot[cA, line width=1.0pt] table[x=f, y=chi_GW241011] {figs/data/e14_scan.dat}; \addlegendentry{GW241011 の主星($\chi_0=0.78$、$10^9$ 年)}
  \addplot[cD, line width=1.0pt] table[x=f, y=chi_M45] {figs/data/e14_scan.dat}; \addlegendentry{$45\Msun$、$\chi_0=0.7$、$10^9$ 年(322 から)}
  \addplot[cC, line width=1.0pt] table[x=f, y=chi_M20] {figs/data/e14_scan.dat}; \addlegendentry{$20\Msun$、$\chi_0=0.3$、$6\times10^9$ 年(322 から)}
  \addplot[only marks, mark=o, mark size=1.6pt, cB] table[x=f, y=chi] {figs/data/e14_ode.dat}; \addlegendentry{三準位の ODE}
\end{axis}
\end{tikzpicture}
